\ssscp{Basic vocabulary: ‘hair’ > ‘leaf of head’}{02-03-01-05}
One might think that basic vocabulary concepts such as ‘hair’
should be fairly stable.
Most languages west of the Blust line retain reflexes of PMP
*bulu ‘body hair, fur, feather’ (data in \qf{02-32} from \citealp{BlustTrussel2020}).

{\wg\tblr{>{\hspace{-\tabcolsep}}l<{\hspace{-\tabcolsep}}>{\hspace{-\tabcolsep}\enspace}lll}
\lsptoprule
\setcounter{tableex}{31}
\tbx{02-32}&	PMP	&	*bulu	&	‘body hair, fur, feather’		\\
\midrule
&	\ili{Kelabit}	&	 \hp{*}buluh	&	‘body hair, feather’		\\
&	\ili{Ngaju} Dayak	&	 \hp{*}bulu	&	‘body hair, fur, feather’		\\
&	\ili{Malay}	&	 \hp{*}bulu	&	‘body hair, fur, feather’		\\
&	Javanese	&	 \hp{*}wulu	&	‘body hair, fur, feather’		\\
&	\ili{Bugis}	&	 \hp{*}bulu	&	‘body hair, fur, feather’		\\
\lspbottomrule
\end{tabular}\mbox{}\\}

\largerpage
This same etymon is also found in many Austronesian languages
of eastern Indonesia and Timor-Leste, albeit following the word order associated with the preposed possessor.
(Data in \qf{02-33} are from Grimes fieldnotes,
Edwards fieldnotes.) Many languages of our target region have
a reflex of *bulu which refers exclusively to ‘body hair,
fur, feather’ and have a different mono-morphemic word for ‘head
hair’.\footnote{
		Many languages have a reflex of PMP *buhək for ‘head hair’.
		Other etyma also occur.}

{\wg\st{5.5pt}
\tblr{>{\hspace{-\tabcolsep}}l<{\hspace{-\tabcolsep}}>{\hspace{-\tabcolsep}\enspace}lr<{\hspace{-\tabcolsep}}>{\hspace{-\tabcolsep}\;}llll}
\lsptoprule
\tbx{02-33}&	\ili{Dela}	&	laŋga 	&	\tbf{fulu-n}	&	‘head hair’	&	\tsc{Timor-Babar} (Rote)		\\
&	\ili{Tii}	&	laŋga 	&	\tbf{ɓulu-n}	&	‘head hair’	&	\tsc{Timor-Babar} (Rote)		\\
&	\ili{Kotos} \ili{Amarasi}	&	ʔnaak 	&	\tbf{funu-n}	&	‘head hair’ 	&	\tsc{Timor-Babar}		\\
&	\ili{Kemak}	&	gala-r	&	\tbf{hulu-n}	&	‘head hair’ 	&	\tsc{\ili{Central} Timor}		\\
&	\ili{Buru}	&	olo-n 	&	\tbf{folo-n}	&	‘head hair’	&	\tsc{Sula-Buru}		\\
&	\ili{Kayeli}	&	olo-ɲ 	&	\tbf{bolo-ɲ}	&	‘head hair’	&	\tsc{\ili{Ambon}-Seram}		\\
\lspbottomrule
\end{tabular}\mbox{}\\}

Yet another pattern begins to be found in the Austronesian languages
in the region as well,
using terms meaning ‘leaf’ -- often reflexes of PMP *dahun ‘leaf’.
\citet[147-148]{Blust1993,DonohueGrimes2008}
and \citet{GrimesCIP} observe that in eastern Indonesia
and Timor-Leste and continuing further eastward into Oceania,
some Austronesian languages have shifted to ‘leaf of head = hair’.
This turns out to be a common Papuan pattern.
Note that in related languages on the same island,
we have both patterns.
Two Rote languages in \qf{02-33} retain reflexes of PMP *bulu,
and two in \qf{02-34} follow the Papuan pattern of ‘head leaf’.
\ili{Termanu} (central Rote) can use either a reflex of *dahun ‘leaf’
or *bulu ‘hair’ for head hair,
with \citet[64]{Jonker1908} giving both phrases as equivalent.
Some of these languages,
such as \ili{Rikou}, retain reflexes of *bulu for ‘animal fur’,
whereas others such as \ili{Dhao} also use a reflex of *dahun ‘leaf’ for ‘animal fur’.
We note that collocation of ‘head leaf’ = ‘hair’ is also found
among Oceanic languages including: \ili{Fijian},
\ili{Tongan}, \ili{Sāmoan}, and \ili{Māori} among others.

{\wg\st{4.7pt}\tblr{>{\hspace{-\tabcolsep}}l<{\hspace{-\tabcolsep}}>{\hspace{-\tabcolsep}\enspace}lr<{\hspace{-\tabcolsep}}>{\hspace{-\tabcolsep}\;}lll}
\lsptoprule
\tbx{02-34} &	\ili{Dhao}	&	\tbf{rəu} 	&	kətu	&	‘\tbf{leaf} of head’ = ‘hair’	&	\tsc{\ili{Bima}-Lembata}		\\
&	\ili{Rikou}	&	laka 	&	\tbf{doo}-ʔ	&	‘head \tbf{leaf}’ = ‘hair’	&	\tsc{Timor-Babar} (Rote)		\\
&	&	laŋa 	&	\tbf{doo}-k	&	‘head \tbf{leaf}’ = ‘hair’	&			\\ \hhline{~}
&		\multirow{-2}{*}{Termanu}	&	laŋa	&	bulu-k	&	‘head hair’	&	\multirow{-2}{*}{\tsc{Timor-Babar} (Rote)}		\\
&	\ili{Welaun}	&	ulu-t 	&	\tbf{loo}-n	&	‘head \tbf{leaf}’ = ‘hair’	&	\tsc{\ili{Central} Timor}		\\
&	\ili{Masiwang}	&	ulin 	&	\tbf{lan}	&	‘head \tbf{leaf}’ = ‘hair’	&	\tsc{STB}		\\
&	\ili{Sekar}	&	unin	&	\tbf{fakin}	&	‘head \tbf{leaf}’ = ‘hair’	&	\tsc{STB}		\\
&	\ili{Arguni}	&	arug 	&	\tbf{nane}	&	‘head \tbf{leaf}’ = ‘hair’	&	\tsc{STB}		\\
\lspbottomrule
\end{tabular}\mbox{}\\}

\largerpage
\ili{As} foreshadowed above, ‘leaf of head’ is a common composition
for ‘head hair’ found in many Papuan languages in eastern Indonesia,
as in \qf{02-35}.

{\wg\st{4.7pt}\tblr{>{\hspace{-\tabcolsep}}l<{\hspace{-\tabcolsep}}>{\hspace{-\tabcolsep}\enspace}lr<{\hspace{-\tabcolsep}}>{\hspace{-\tabcolsep}\;}lll}
\lsptoprule
\tbx{02-35}&	Pura	&	oŋ 	&	\tbf{vaa}	&	‘head \tbf{leaf}’ = ‘hair’	&	TAP		\\
&	\ili{Teiwa}	&	g-oʔon 	&	\tbf{waʔ}	&	‘\tsc{3sg}-head \tbf{leaf}’ = ‘hair’	&	TAP		\\
&	\ili{Makasae}	&	daʔe (gi) 	&	\tbf{asa}	&	‘head (\tsc{poss}) \tbf{leaf}’ = ‘hair’	&	TAP		\\
&	\ili{Iha}	&	kanda 	&	\tbf{den}	&	‘head \tbf{leaf}’ = ‘hair’	&	\tsc{West Bomberai}		\\
&	\ili{Mbaham}	&	kanda 	&	\tbf{tawe}	&	‘head \tbf{leaf}’ = ‘hair’	&	\tsc{West Bomberai}		\\
\lspbottomrule
\end{tabular}\mbox{}\\}

%\begin{figure}
	%\includegraphics[width=\columnwidth]{head-leaf.pdf}
	%\caption[Distribution of ‘leaf of head’ = ‘hair’]
	%{Distribution of ‘leaf of head’ = ‘hair’}
	%\label{fig:02-18}
%\end{figure}

\begin{figure}
	\includegraphics[width=\columnwidth]{head-leaf.pdf}
	\caption[Distribution of ‘leaf of head’ = ‘hair’]
	{Distribution of ‘leaf of head’ = ‘hair’\footnotemark}
	\label{fig:02-18}
\end{figure}
	\footnotetext{
	Red circles/squares in \frf{02-18} indicate languages where ‘head
	hair’ is transparently, synchronically analysable as composed of ‘head’ + ‘leaf’.
	Yellow circles/squares, (*head) + *leaf are for languages in
	which the word for ‘head hair’ contains a term which historically
	meant ‘leaf’ but no longer has this synchronic meaning.
	Examples of the latter include: Northwest \ili{Mambae} (Hatulia) \it{ulu
	laun} < *qulu ‘head’ + *dahun ‘leaf’ but \it{glutan} ‘head’ \it{nora}
	‘leaf’, Palu{\Q}e with \it{lolo} ‘head hair’ (cognates of which in related
	languages mean ‘leaf’; e.g.
	\tsc{Western Lamaholot} \it{lolon} ‘leaf’) but Palu{\Q}e \it{ʋunu}
	‘leaf’, or \ili{Moor} \it{ʔuranu} ‘head hair’ with \it{ʔranu} ‘leaf’ but \it{βarota}
	‘head’.}

So following the logic spelled out by \citet[14--16]{Klamer2019},
it is reasonable to infer that the pattern ‘head leaf’ = ‘hair’
in Austronesian languages \ili{east} of the Blust line is triggered
by contact with Papuan languages.
We note that its distribution is sporadic.
It does not replace PMP *bulu,
and does not cover all languages in any subgroup.
These factors mean we cannot use it for linguistic classification.
But it is also geographically widespread \ili{east} of the Blust line
region persisting as far \ili{east} as Polynesia.
It is a “contaminant” from the perspective of understanding Austronesian
dispersal, while at the same time being a very valuable trace of the intensity
of contact with Papuan languages,
even very far from where they are still spoken.
The distribution of known reflexes of ‘leaf of head’ = ‘hair’
in Papuan and Austronesian languages in the Wallacean region is shown in
\frf{02-18}.
